\documentclass[oneside,12pt]{book}
\usepackage[utf8]{inputenc}
\usepackage{standalone}
\usepackage[margin=2.5cm]{geometry}
\usepackage{graphicx}
\usepackage{hyperref}
\usepackage{lmodern}
\usepackage{float}
\usepackage[spanish]{babel}
\usepackage{csquotes}
\usepackage[figurename=Ilustración]{caption}
\usepackage[numbers]{natbib}
\usepackage{ragged2e}
\usepackage{array}
\usepackage{multirow}
\usepackage{parskip}
\usepackage{listings}
\usepackage{epigraph}
\usepackage{lipsum} % Texto de relleno en plantilla

\usepackage{comment} %Uso de comentarios multilínea

%\usepackage{fancyhdr}
\usepackage{geometry}
\usepackage{color}\usepackage{graphicx,xcolor}

\renewcommand{\lstlistingname}{Algoritmo}
\renewcommand{\lstlistlistingname}{Índice de~\lstlistingname s}
\lstset{
        basicstyle=\fontsize{7}{11}\ttfamily,
        language=Java,
        captionpos=b,
}

\usepackage{titlesec}
\setcounter{secnumdepth}{3}
\titleformat{\chapter}[display]
{\normalfont\huge\bfseries}{\chaptertitlename\ \thechapter}{20pt}{\Huge}

\urlstyle{same}
\title{Galería de NFT's en la Web3}
\author{Cristóbal José Jiménez Gómez}
\date{Curso 2022 - 2023}
\renewcommand{\contentsname}{Contenido}
\renewcommand{\figurename}{Ilustración}
\usepackage{pifont}
\renewcommand\labelitemi{\ding{52}}

\begin{document}
\begin{titlepage}
    \begin{center}
    \vspace*{0.5in}
        \begin{figure}[htb]
            \begin{center}
            \includegraphics[width=15cm]{Ilustraciones/LogoEII.jpg}
            \end{center}
        \end{figure}

        \vspace*{0.15in}
        \vspace*{0.2in}

        \noindent\hfil\rule{11cm}{0.2mm}\hfil\\
        \vspace*{0.1in}
        \begin{Huge}
            \textbf{Galería de NFT's en Web3} \\
        \end{Huge}

        \vspace*{0.3in}

        \begin{large}
            TITULACIÓN: Grado en Ingeniería Informática \\
            \vspace*{0.1in}
            AUTOR: Cristóbal José Jiménez Gómez \\
        \end{large}
        \vspace*{0.3in}
        
        \noindent\hfil\rule{11cm}{0.2mm}\hfil\\

        \vspace*{0.1in}

        \begin{large}
            TUTORIZADO POR: \\
            María Dolores Afonso Suárez \\
        \end{large}
        
    \vspace*{0.3in}

    Fecha: 05/2023
    \end{center}
\end{titlepage}


%\newgeometry{margin=1.5cm}

\thispagestyle{empty}

%\lhead{ 
\begin{tabular}{p{15cm}p{5cm}}
  \includegraphics{Ilustraciones/LogoEII.jpg}   &  TFT04\\
\end{tabular}
%}

\vspace{1em}
\fboxrule=2pt
\begin{center}
\fcolorbox{gray}{gray!10}{\parbox{.8 \textwidth}{ {\large SOLICITUD DE DEFENSA DE TRABAJO DE FIN DE TÍTULO}}}
\end{center}

\vspace{1em}
\justify
D./Dª Nombre\_Estudiante, autor del trabajo Título\_Trabajo, correspondiente a la titulación Titulación\_Estudiante, en colaboración con la empresa/proyecto Nombre\_Empresa (en su caso)

\vspace{1em}
SOLICITA

\vspace{1em}
que se inicie el procedimiento de defensa del mismo, para lo que se adjunta la documentación requerida, haciendo constar que 

[ ] se autoriza / [ ] no se autoriza la grabación en audio de la exposición y turno de preguntas.

\vspace{1em}
Asimismo, con respecto al registro de la propiedad intelectual/industrial del TFT, declara que:

[ ] Se ha iniciado o hay intención de iniciarlo (defensa no pública).

[ ] No está previsto.

\vspace{1em}
Y para que así conste firma la presente. (fecha en firma electrónica)

\begin{center}
%En Las Palmas de Gran Canaria, a día de mes de año

\vspace{1em}
El/La estudiante

\vspace{3em}
Fdo.------------------------
\end{center}

\begin{center}
\fcolorbox{gray}{white!10}{\parbox{.9\textwidth}{
{\tiny A rellenar y firmar \textbf{obligatoriamente} por el/la/los/las tutores}

En relación con la presente solicitud, se informa:{\tiny (firmar donde corresponda)}

\vspace{1em}
\begin{center}
\begin{tabular}{|p{\dimexpr.45\linewidth}|p{\dimexpr.45\linewidth}|}
%\begin{tabular}{|l|r|}
    \hline
   & \\
  Positivamente {\tiny (en  caso  de  detección  de  copia, esta firma quedará invalidada)}  & Negativamente {\tiny (justificación en TFT05)} \\
   & \\
   & \\
   & \\
   & \\
   & \\
   & \\
  Fdo.------------------------ & Fdo.------------------------  \\
     & \\

    \hline
\end{tabular} 

\end{center}

\vspace{1em}

}
}

\vspace{1em}

%\cfoot{
DIRECTORA DE LA ESCUELA DE INGENIERÍA INFORMÁTICA
%}

\end{center}

\restoregeometry

\clearpage

\newpage
\pagenumbering{arabic}
{\Large{\textbf{Agradecimientos}}}

\begin{flushright}
    {\setlength{\parskip}{8mm}
        \large{
            \textit{A los gnomos de jardín.}
        }
    }
    
\end{flushright}

\clearpage
{\setlength{\parskip}{8mm}
{\huge{\textbf{Resumen}}}

    {\large{ 
        Desarrollo de una Aplicación Web en el marco de la Web3(blockchain), que gestiona una galería de 
        Non Fungible Tokens (NFT’s). Ofrece la gestión de contenidos a distintos tipos de usuarios cuyos niveles de
        interacción varían desde los usuarios sin registrar, a administradores. En esa escala de accesos los
        casos de uso establecerán el nivel de interacción de cada perfil. Las funcionalidades implementadas
        permitirán la gestión a distintos niveles de estos NFT’s, desde la visualización hasta la gestión de cada
        cartera. Este último concepto es el definido para almacenarlos y, en caso de ser necesario comerciar
        (compraventa o intercambio) con ellos.
    }}

    \textbf{Palabras claves:} Web3, Web3.0, Angular, Firebase, Vercel, Express
    

%\newpage

{\huge{\textbf{Abstract}}}

    {\large{
        Development of a Web Application within the framework of Web3 (blockchain), which manages a gallery of
        Non Fungible Tokens (NFT's). It offers content management to different types of users whose levels of
        Interaction range from unregistered users to administrators. At this access scale, the
        use cases will establish the level of interaction of each profile. The functionalities implemented
        will allow the management at different levels of these NFTs, from the visualization to the management of each
        briefcase. This last concept is the one defined to store them and, in case it is necessary to trade
        (purchase or exchange) with them.
    }}

    \textbf{Key Words:} Web3, Web3.0, Angular, Firebase, Vercel, Express
} 
\tableofcontents
\listoffigures
\listoftables
\lstlistoflistings

\clearpage
\pagenumbering{arabic}

\chapter{Introducción}
\section{Motivación}
La razón principal que se encontró para la elaboración de este trabajo ha sido 
el auge de la Blockchain durante estos últimos años. "Blockchain es un libro 
mayor compartido e inalterable que facilita el proceso de registro de 
transacciones y de seguimiento de activos en una red de negocios. Un activo 
puede ser tangible o intangible. Prácticamente cualquier cosa de valor, puede 
rastrearse y comercializarse en una red de blockchain, reduciendo así el 
riesgo y los costes para todos los involucrados"\cite{ibmWeb3}. 

Esta tecnología ofrece la posibilidad de cosntruir aplicaciones web que garanticen
la privacidad y seguridad de los datos de un usuario. La Web 3 ofrece transparencia
y trazabilidad a la hora de la interacción entre usuarios. Además, permite la 
creación o conversión de negocios ya afianzados a este entorno.

Además, se afianza una necesidad que encontraba yo, personalmente, a la hora de
interactuar entre distintas carteras, que es la necesidad de tener todas en una 
misma ubicación.

Se debe constar que ya existen páginas afianzadas para el fin de mostrar y 
compra-venta de NFT's. Este trabajo busca el aprendizaje tanto del desarrollo 
de una aplicación web completa, como de aprender a realizar las llamadas y 
conexiones para recoger y mostrar datos de NFT's\cite{bbcNFTs} y 
criptocarteras\cite{moralisWallets}.

\newpage
\section{Objetivos}
El objetivo principal del proyecto se centra en el desarrollo de un sitio web 
que utiliza el nuevo concepto de descentralización dentro del marco Web3. 
El resultado ofrece como funcionalidad principal la búsqueda de NFT’s y el 
almacenamiento de distintas carteras de criptomonedas. Para ello se plantea 
un conjunto de objetivos generales:
\begin{enumerate}
    \item Estudiar las características y evolución de la Web 3.0 y su relación con la Web3.
    \item Realizar un estudio de las tecnologías y herramientas a emplear.
    \item Establecer criterios para su selección (tanto las tecnologías como de las herramientas)
    \item Seguir una metodología de desarrollo, las pruebas y documentación del proyecto.
\end{enumerate}




\section{Estructura de la memoria}
La memoria está dividida en 10 capítulos diferentes. En ellos se recoge toda
la información, desde el inicio del proyecto hasta las conclusiones. Así, 
la información recogida en cada capítulo se describe tal que:

\-\hspace{1cm} El primer capítulo consta de la introducción al proyecto. En este se 
contextualiza el trabajo, se explica la motivación detrás de la elaboración 
del proyecto, se dan a conocer los objetivos a cumplir y se explica la 
estructura del documento.

\-\hspace{1cm} El segundo capítulo lista y justifica las competencias de la titulación
abarcadas durante el desarrollo del trbajo mediante las teras y actividades realizadas.

\-\hspace{1cm} El tercer capítulo desarrolla el marco teórico, donde se explican y 
analizan algunos términos relevantes a este trabajo. Es el punto de partida, 
donde se realiza una investifación que contextualiza y justifica las cacciones
que se toman en el trabajo.

\-\hspace{1cm} El cuarto capítulo abarca el estado del arte, que explica la 
actualidad del desarrollo Web junto con el desarrollo dentro de la Web3.0 y 
la popularidad entre los distintos frameworks\footnote{Un framework es una 
herramienta de programación que te permite desarrollar software proporcionando 
una estructura con componentes integrados que sirven de base para construir 
proyectos nuevos\cite{bootcampFramework}}

\-\hspace{1cm} El quinto capítulo contiene los recursos y las tecnologías 
utilizadas para el desarrollo del trabajo. En este se definen y justifican.

\-\hspace{1cm} El sexto capítulo se explica la metodología llevada a cabo 
para la realización del trabajo, se dan a conocer als diferentes fases y 
tareas que han existido durante su desarrollo.

\-\hspace{1cm} El séptimo capítulo se incluye un análisis de los requisistos y 
el diseño de la página que se ha elaborado para el desarrollo.

\-\hspace{1cm} El octavo capítulo se explica todo lo relacionado con el 
desarrollo de la página web.

\-\hspace{1cm} El noveno capítulo se encuentra la evaluación de la calidad de la página, 
obtenidos a través de diversos programas o extensiones.

\-\hspace{1cm} El décimo capítulo se incluye las conclusiones que se pueden obtener 
de este proyecto, tanto personales como objetivas, dando además una propuesta de mejora 
para un futuro desarrollo.

\-\hspace{1cm} Finalmente se incluyen los anexos, los cuales también contienen las 
referencias bibliográficas



%
\chapter{Estado actual y objetivos iniciales}
Estado del arte to rechulón
CAPÍTULO OBLIGATORIO

Descripción del estado actual en la temática concreta del trabajo y desglose de los objetivos inicialmente previstos en el TFT01.

\chapter{Competencias específicas y aportaciones del trabajo}
%\input{Competencias.tex}
CAPÍTULO OBLIGATORIO

Indicar, sólo para las competencias específicas relacionadas de forma más directa con el trabajo desarrollado, cómo se han cubierto con este TFT.

Justificar qué es lo que este TFT aporta a nuestro entorno socio-económico, técnico o científico.

\chapter{Desarrollo}
%\input{Desarrollo.tex}
CAPÍTULO OBLIGATORIO

Metodología aplicada y desarrollo del trabajo en sus distintas fases, decisiones de diseño, herramientas empleadas.

Si el trabajo incluye software desarrollado, deberán seleccionarse las secciones más relevantes del mismo y comentarlas en la memoria.

Ajuste a la planificación inicialmente prevista.

Modificación en los objetivos planteados.

\chapter{Conclusiones y trabajo futuro}
%\input{Conclusiones.tex}
CAPÍTULO OBLIGATORIO

Resultados,  grado  de  consecución  de  los  objetivos,  posibles extensiones

\chapter{OTROS}

El documento debe terminar con las referencias bibliográficas (SECCIÓN OBLIGATORIA). Después pueden incluirse anexos 

CAPÍTULOS OPCIONALES

Dependiendo del tipo de trabajo, se podrían incluir capítulos adicionales como los siguientes:

\begin{description}
\item{REQUISITOS}
\item{DISEÑO}
\item{MANUAL  DE  USUARIO  Y  SOFTWARE}  deberán  incluirse  obligatoriamente  en  la  memoria  los extractos más relevantes   del código desarrollado. Siempre que sea posible, deberá   proporcionarse acceso a un repositorio software.
\item{NORMATIVA Y LEGISLACIÓN} incluir la legislación vigente que afecte al TFT (ley de protección de datos, leyes sobre seguridad, ...)
\item{ASPECTOS ECONÓMICOS Y TEMPORALES}
\end{description}



\bibliography{ref}
\bibliographystyle{apalike}

\end{document}